\chapter{Observabilité — Prometheus et Grafana}
\label{ch:monitoring}

\section{Objectifs}

L'observabilité vise à répondre, en exploitation :
\begin{itemize}
  \item l'API reçoit-elle du trafic ? à quel rythme ?
  \item quelle est la latence (p50, p95, p99) ?
  \item y a-t-il des erreurs 5xx ?
  \item quelles routes sont les plus sollicitées ?
\end{itemize}

StockFlow adopte la stack \textbf{Prometheus} (collecte) + \textbf{Grafana} (visualisation), standard dans l'écosystème Kubernetes.

\section{Instrumentation applicative}

Le package \texttt{prometheus-fastapi-instrumentator} expose automatiquement sur \texttt{GET /metrics} :
\begin{itemize}
  \item compteur \texttt{http\_requests\_total} (labels method, handler, status) ;
  \item histogrammes de durée de requête ;
  \item métriques en vol (\texttt{inprogress}).
\end{itemize}

Aucune modification manuelle par route n'est nécessaire pour les métriques HTTP de base.

\section{kube-prometheus-stack}

Terraform installe le chart communautaire \texttt{kube-prometheus-stack} (version 65.5.0) avec des values dans \texttt{monitoring/kube-prometheus-values.yaml} :

\begin{itemize}
  \item \textbf{Grafana} activé, mot de passe admin via \texttt{set\_sensitive} Terraform ;
  \item \textbf{sidecar dashboards} — surveille les ConfigMaps labellisées \texttt{grafana\_dashboard: "1"} dans tous les namespaces ;
  \item \textbf{Prometheus} — rétention 1 jour, ressources limitées pour VPS ;
  \item composants control-plane désactivés pour alléger le cluster ;
  \item \textbf{Grafana} exposé sur \texttt{grafana.pfa.elyes.dev} avec TLS (ingress + cert-manager).
\end{itemize}

\section{ServiceMonitor}

Le CRD \texttt{ServiceMonitor} (Prometheus Operator) décrit comment scraper une cible. Le chart \texttt{pfa-stock} crée :

\begin{lstlisting}[style=yaml,caption={Extrait ServiceMonitor (conceptuel)}]
spec:
  selector:
    matchLabels:
      app: api
  endpoints:
    - port: http
      path: /metrics
      interval: 30s
\end{lstlisting}

Le label \texttt{release: kube-prometheus-stack} associe ce monitor à l'instance Prometheus déployée par le chart homonyme.

\section{Flux de métriques}

\begin{figure}[htbp]
  \centering
  \begin{tikzpicture}[node distance=1.1cm]
    \node[block, svc] (pod) {Pod API\\/metrics port 8000};
    \node[block, infra, right=1.6cm of pod] (sm) {ServiceMonitor};
    \node[block, infra, right=1.6cm of sm] (prom) {Prometheus};
    \node[block, cloud, below=1.2cm of prom] (graf) {Grafana};
    \draw[line] (pod) -- (sm);
    \draw[line] (sm) -- (prom);
    \draw[line] (prom) -- (graf);
  \end{tikzpicture}
  \caption{Chemin des métriques}
  \label{fig:metrics-flow}
\end{figure}

\section{Dashboard Grafana}

Le fichier \texttt{monitoring/dashboards/pfa-stock-api.json} définit le tableau de bord \emph{PFA Stock — API FastAPI} avec panneaux :
\begin{itemize}
  \item débit de requêtes ;
  \item erreurs 5xx ;
  \item latences p50 / p95 / p99 ;
  \item répartition par route et par code HTTP ;
  \item requêtes en cours.
\end{itemize}

Terraform injecte ce JSON en ConfigMap ; le sidecar Grafana le charge automatiquement — le dashboard est \textbf{versionné comme le code}.

\section{Accès opérationnel}

Grafana est accessible en production sur \texttt{https://grafana.pfa.elyes.dev} (identifiants définis au \texttt{terraform apply}).

\section{Mode développement}

En \texttt{make dev}, Prometheus/Grafana ne sont pas déployés par défaut. Les métriques restent accessibles sur \texttt{http://localhost:8000/metrics} pour curl ou un Prometheus local ad hoc.

\section{Alerting}

Alertmanager est inclus dans kube-prometheus-stack mais les règles d'alerte custom (stock critique, API down prolongé) ne sont pas configurées dans le périmètre actuel — extension naturelle pour une phase ultérieure.
